%==============================================================================
%  TEMPLATE CHUẨN CHO MỌI FILE .TEX TRONG REPO NÀY
%  Phong cách: đề/lời giải kiểu Harvard-MIT Math Tournament (HMMT)
%  - Chỉ trắng và đen: không màu, không khung, không tcolorbox, không đổ bóng.
%  - Tiêu đề gọn ở giữa trang, không trang bìa, không \maketitle.
%  - Đề bài chạy dòng ("1." / "Bài 2.1." in đậm), lời giải run-in "Lời giải:".
%  - Đáp số cuối cùng đóng khung bằng \boxed{...}.
%  Biên dịch: xelatex (bắt buộc, vì dùng fontspec + polyglossia + unicode-math)
%==============================================================================
\documentclass[11pt,a4paper]{article}

% --- Ngôn ngữ và font: Latin Modern (Computer Modern) như bản HMMT gốc -------
\usepackage{fontspec}
\usepackage{polyglossia}
\setdefaultlanguage{vietnamese}
\setotherlanguage{english}
\setmainfont{lmroman10-regular.otf}[
  BoldFont       = lmroman10-bold.otf,
  ItalicFont     = lmroman10-italic.otf,
  BoldItalicFont = lmroman10-bolditalic.otf
]

% --- Toán học ---------------------------------------------------------------
\usepackage{amsmath, amssymb, amsthm}
\usepackage{unicode-math}
\setmathfont{latinmodern-math.otf}
\allowdisplaybreaks

% --- Khổ trang: lề rộng, khối chữ gọn vào giữa ------------------------------
\usepackage[a4paper, top=3.0cm, bottom=3.0cm, left=3.2cm, right=3.2cm]{geometry}
\usepackage{enumitem}
\setlist{itemsep=2pt, topsep=2pt, parsep=0pt, partopsep=0pt, leftmargin=1.6em}

% Đoạn văn kiểu HMMT: thụt đầu dòng, KHÔNG dùng parskip (không chèn dòng trống)
\setlength{\parindent}{1.5em}
\setlength{\parskip}{0pt}
\setlength{\abovedisplayskip}{5pt plus 1pt minus 2pt}
\setlength{\belowdisplayskip}{5pt plus 1pt minus 2pt}
\setlength{\abovedisplayshortskip}{2pt plus 1pt}
\setlength{\belowdisplayshortskip}{2pt plus 1pt}

% Chỉ số trang canh giữa chân trang, không running header
\pagestyle{plain}

% --- Thông tin tài liệu (sửa 4 dòng này cho mỗi bài) ------------------------
\newcommand{\coursename}{High Dimensional Statistics}
\newcommand{\docdate}{\today}
\newcommand{\doctitle}{Exercise 1: Lời giải}
\newcommand{\studentname}{Nguyễn Quang}

% --- Đề bài / lời giải / định lý -------------------------------------------
\newcounter{prob}
\makeatletter
% \begin{problem}            -> tự đánh số 1., 2., 3., ...
% \begin{problem}[Bài 2.1.]  -> dùng nhãn tự đặt
\newenvironment{problem}[1][]{%
  \par\medskip\noindent
  \def\@tempa{#1}%
  \ifx\@tempa\@empty\refstepcounter{prob}\textbf{\theprob.}\else\textbf{#1}\fi
  \ \ignorespaces
}{\par}
\makeatother

% Lời giải: run-in ngay dưới đề bài, không xuống dòng thừa, không ký hiệu ∎
\newenvironment{solution}{%
  \par\noindent\textbf{Lời giải:}\ \ignorespaces
}{\par}

% Nhắc lại lý thuyết: kiểu amsthm run-in, chữ đứng, không khung
\theoremstyle{definition}
\newtheorem{thm}{Định lý}
\newtheorem{dfn}{Định nghĩa}

% Đáp số cuối cùng luôn đóng khung
\newcommand{\answer}[1]{\boxed{#1}}

\begin{document}

%------------------------------ Tiêu đề gọn ---------------------------------
\begin{center}
  {\bfseries\coursename}\\
  \docdate\\
  \doctitle: \textbf{\studentname}
\end{center}
\vspace{0.5em}

%------------------------------ Nội dung ------------------------------------
\begin{problem}[Bài 2.1.]
Cho $\mathbf{X} = (X_1, X_2, X_3)' \sim \mathcal{N}_3(\boldsymbol{\mu}, \boldsymbol{\Sigma})$ với
\[
\boldsymbol{\mu} = \begin{pmatrix} 1 \\ -1 \\ 2 \end{pmatrix}, \qquad
\boldsymbol{\Sigma} = \begin{pmatrix} 4 & 0 & -1 \\ 0 & 5 & 0 \\ -1 & 0 & 2 \end{pmatrix}.
\]
Xác định phân phối của $(X_1, X_2)'$.
\end{problem}

\begin{solution}
Đặt $\mathbf{A} = \begin{pmatrix} 1 & 0 & 0 \\ 0 & 1 & 0\end{pmatrix}$ thì
$(X_1, X_2)' = \mathbf{A}\mathbf{X} \sim \mathcal{N}_2(\mathbf{A}\boldsymbol{\mu}, \mathbf{A}\boldsymbol{\Sigma}\mathbf{A}')$,
do đó
\[
\begin{pmatrix} X_1 \\ X_2 \end{pmatrix} \sim
\answer{\mathcal{N}_2\!\left(\begin{pmatrix} 1 \\ -1 \end{pmatrix},
\begin{pmatrix} 4 & 0 \\ 0 & 5 \end{pmatrix}\right)}.
\]
Vì hiệp phương sai chéo bằng $0$ và hai biến chuẩn đồng thời nên $X_1 \perp X_2$.
\end{solution}

\begin{problem}
Ví dụ đề bài tự đánh số: nêu điều kiện để $\mathbf{U}$ và $\mathbf{V}$ chuẩn đồng thời là độc lập.
\end{problem}

\begin{thm}[Tính độc lập trong phân phối chuẩn nhiều chiều]
Nếu $\mathbf{U}$ và $\mathbf{V}$ có phân phối chuẩn đồng thời thì
$\mathbf{U} \perp \mathbf{V} \iff \operatorname{Cov}(\mathbf{U}, \mathbf{V}) = \mathbf{O}$.
\end{thm}

\begin{solution}
Hệ quả trực tiếp của định lý trên: chỉ cần kiểm tra khối hiệp phương sai chéo có bằng $\mathbf{O}$ hay không.
\end{solution}

\end{document}
